\section{The ChainBench-ADD Benchmark}
\label{sec:benchmark}

\subsection{Overview and Terminology}
ChainBench-ADD organizes post-generation delivery into families, templates, variants, and realized chains. This hierarchy supports scenario-, route-, local-, and path-level analysis under matched-parent control.

Throughout this section, a delivery operator denotes an atomic or composite transform such as codec compression, re-encoding, packet loss, room reverberation, or additive noise. A template is an ordered operator blueprint within one delivery family, and a variant is one parameterized instantiation. A delivery chain is the realized ordered path applied to one parent sample. A delivery family is the top-level scenario category that groups related templates, including direct, platform-like, telephony, simulated replay, and hybrid delivery.

\subsection{Delivery Families and Chain Realization}
ChainBench-ADD defines five delivery families that reflect common post-generation routes. Direct is the identity control. Platform-like approximates online redistribution and recompression. Telephony models bandwidth-limited call transmission with telephony codecs, packet impairment, and compact end-to-end call simulation. Simulated replay models playback and recapture in physical spaces through room response and environmental noise. Hybrid combines communication and acoustic effects within a single path. These families are best understood as scenario categories rather than disjoint operator inventories: different families may share operators but combine them in different ordered patterns and parameter regimes.

These families are built from eight reusable delivery operators: resampling, band-limiting, codec compression, re-encoding, packet loss, additive noise, room impulse response (RIR), and call-path simulation. Many operators are shared across multiple families. For example, codec-like effects appear in both platform-like and telephony delivery, while RIR appears in both simulated replay and hybrid delivery. A template specifies an ordered operator blueprint within one delivery family, a variant instantiates that blueprint with concrete parameter choices, and a realized chain is the waveform-domain rendering of that ordered path for one parent sample. Table~\ref{tab:ChainBench-ADD-operators} summarizes the operator inventory and representative parameter ranges of each operator. Full parameter pools are deferred to the appendix.

\begin{table}[t]
\caption{Delivery operators in ChainBench-ADD.}
\label{tab:ChainBench-ADD-operators}
\vspace{-0.5em}
\centering
\footnotesize
\setlength{\tabcolsep}{2pt}
\renewcommand{\arraystretch}{1.30}
\begin{tabularx}{\linewidth}{@{}%
    >{\raggedright\arraybackslash}p{0.17\linewidth}%
    @{\hspace{2pt}}%
    >{\raggedright\arraybackslash}X%
    @{\hspace{2pt}}%
    >{\centering\arraybackslash}p{0.08\linewidth}%
    @{}}
\toprule
\textbf{Operator} & \textbf{Representative settings} & \textbf{Fam.} \\
\midrule
Resample 
& 16$\rightarrow$8, 16$\rightarrow$24, 16$\rightarrow$8$\rightarrow$16, 
16$\rightarrow$24$\rightarrow$16, 16$\rightarrow$32$\rightarrow$16 kHz 
& P/T/R/H \\

Band-limit 
& Narrowband, wideband 
& T/H \\

Codec 
& AAC (24/32/48 kbps), Opus (16/24/32 kbps), GSM, $\mu$-law PCM 
& P/T/H \\

Re-encode 
& Same/cross-codec AAC/Opus recompression (24/32 kbps) 
& P/R/H \\

Packet loss 
& 1/3/5/10\% loss; burst 2/3/5; repeat-fade, interpolation, or noise-fill concealment 
& T/H \\

Noise 
& White, pink, brown, babble, hiss, hum; 30/20/15/10 dB SNR 
& R \\

RIR 
& Small/medium/large rooms; RT60 0.2/0.4/0.6/0.8 s; source--microphone distance 0.5/1/2/3 m 
& R/H \\

Call-path 
& Joint channel filtering, codec, packet loss, jitter (0/8/16 ms), and AGC 
& T \\
\bottomrule
\end{tabularx}

\medskip
\raggedright
\footnotesize \textbf{Fam.} P = platform-like, T = telephony, R = simulated replay, H = hybrid.
\end{table}

All operators are implemented as waveform-domain transforms rather than symbolic tags. Codec and re-encode effects use actual encode--decode round trips, packet loss includes burstiness and concealment, RIR is applied with Pyroomacoustics~\cite{scheibler2018pyroomacoustics}, and call-path simulation jointly models channel filtering, codec compression, burst loss, jitter-buffer effects, and AGC. After rendering, all outputs are standardized to 16\,kHz mono PCM WAV.

This hierarchy supports analysis at multiple levels, including scenario-level robustness, generalization to unseen route patterns, operator-level sensitivity, and path-wise robustness under accumulated edits. The current release contains 33 templates in total (1 direct, 6 platform-like, 10 telephony, 7 simulated replay, and 9 hybrid), which realize 26 unique ordered operator sequences and are rendered into sample-specific realized chains. Figure~\ref{fig:chainbench-overview} summarizes this structure.

\begin{figure*}[t]
  \centering
  \captionsetup[subfigure]{justification=centering}
  \begin{subfigure}[t]{0.28\textwidth}
    \centering
    \includegraphics[width=\linewidth]{figures/chainbench_add_a_family_v5.pdf}
    \caption{Delivery-family composition.}
    \label{fig:chainbench-family}
  \end{subfigure}\hfill
  \begin{subfigure}[t]{0.34\textwidth}
    \centering
    \includegraphics[width=\linewidth]{figures/chainbench_add_b_family_operator_v5.pdf}
    \caption{Delivery-family/operator map.}
    \label{fig:chainbench-family-operator}
  \end{subfigure}\hfill
  \begin{subfigure}[t]{0.30\textwidth}
    \centering
    \includegraphics[width=\linewidth]{figures/chainbench_add_c_transition_v5.pdf}
    \caption{Operator transitions.}
    \label{fig:chainbench-transition}
  \end{subfigure}
  \caption{Composition and structure of ChainBench-ADD delivery chains. (a) Sample counts by delivery family, split into bona fide and spoof. (b) Template-level delivery-family/operator map, where each cell shows how many templates in a family contain the operator. (c) Adjacent-operator transition matrix induced by the template inventory.}
  \label{fig:chainbench-overview}
\end{figure*}

\subsection{Data Construction}
ChainBench-ADD is built with a staged pipeline that keeps transcript and speaker identity fixed while introducing spoofing and delivery effects in separate steps. This design follows the benchmark goal: to measure robustness to post-generation delivery artifacts rather than differences in content, speaker, or source formatting. We therefore first construct a clean parent pool, then generate spoof counterparts, and only afterward render delivery chains and export protocol-aware metadata. This ordering ensures that delivery is evaluated as a downstream factor rather than being conflated with generator choice, transcript content, or speaker identity.

We curate bona fide speech from Common Voice~24.0 for English~\cite{ardila2020commonvoicemassivelymultilingualspeech} and AISHELL-3~\cite{shi2020aishell} for Mandarin Chinese. Candidate utterances are screened at both the text and audio levels, removing unsuitable transcripts and low-quality recordings. Selection is speaker-centric: for each accepted speaker, a fixed number of utterances is retained after ranking candidates by properties such as usable duration, speech activity, and stable loudness. We then assign benchmark speaker IDs and create train/dev/test splits at the speaker level. The retained utterances are converted into clean masters by trimming silence, downmixing to mono, resampling to 16\,kHz, normalizing loudness, and saving as 16-bit PCM WAV, yielding 18,703 bona fide clean parents.

For each bona fide clean parent, we synthesize two spoof clean parents using two generators sampled from six contemporary systems: Qwen3-TTS~\cite{Qwen3-TTS}, CosyVoice3~\cite{du2025cosyvoice3}, Spark-TTS~\cite{wang2025sparktts}, F5-TTS~\cite{chen-etal-2024-f5tts}, IndexTTS2~\cite{zhou2025indextts2}, and VoxCPM~\cite{voxcpm2025}. The generator set is chosen to span multiple synthesis families, including multilingual instruction TTS, zero-shot cloning TTS, expressive controllable TTS, and architecture-diverse TTS. To preserve identity while avoiding trivial prompt leakage, the target transcript is paired with a different clean utterance from the same speaker as prompt/reference audio whenever possible. After validation, this stage yields 37,110 spoof clean parents and 55,813 clean parents in total.

Delivery chains are rendered only after the clean pool is fixed, keeping delivery orthogonal to the bona fide/spoof label. For each parent, we generate one direct control and sampled variants from the non-direct families by instantiating ordered waveform-domain operator chains. The pipeline also preserves paired constructions for controlled local comparisons, including order-swap protocols. Final packaging validates audio and exports metadata that links each sample to its parent, family, template, variant, and operator sequence. This explicit lineage is what later enables matched local interventions and path-wise robustness analysis under shared parent context.

\subsection{Benchmark Tasks and Protocols}
ChainBench-ADD is released as a metadata-driven benchmark for delivery-aware audio deepfake detection. From the shared metadata, we instantiate five tasks: in-chain detection, operator substitution, parameter perturbation, order swap, and delivery robustness.

The base task is in-chain detection, which performs bona fide-vs.-spoof classification separately within each delivery family. The three intervention tasks are constructed under matched parent, family, and label control. In operator substitution, paired samples have the same parent utterance, delivery family, ground-truth label, and chain length, but differ only in the operator identity at one or more positions. In parameter perturbation, paired samples differ at exactly one operator position while keeping the operator identity fixed and changing only its parameterization. In order swap, paired samples preserve the same operators and label but differ by one adjacent minimal swap in operator order. Finally, delivery robustness groups samples into lineage graphs defined by shared parent utterance, delivery family, and label, and evaluates whether correct decisions persist as additional observed delivery edits accumulate along the chain. Combined with an extra leave-one-template-out evaluation, these protocols test family-level robustness, transfer to unseen templates, sensitivity to controlled structural edits, and stability along realistic delivery paths.

\subsection{Dataset Statistics, Metadata, and Validation}
ChainBench-ADD contains 941,201 waveforms derived from 55,813 parents and 448 speakers across English and Mandarin Chinese, including 315,573 bona fide and 625,628 spoof samples. The speaker-disjoint train/dev/test split contains 663,363/136,027/141,811 samples from 314/67/67 speakers, respectively. At the delivery level, the release covers five delivery families, 33 templates, and 26 ordered operator sequences. At the protocol level, it further provides 182,409 operator-substitution groups, 53,234 parameter-perturbation groups, 108,245 order-swap groups, and 27,204 lineage/path groups.

The released metadata are lightweight and protocol-aware. In addition to identifiers, transcript, language, speaker, and provenance, they include split label, generator family and model, delivery family, template and variant identifiers, ordered operator signatures, and the structural annotations needed to reconstruct the benchmark tasks. Before export, each waveform is rechecked against target format and duration constraints, and the final package records split-integrity checks, duplicate-identifier and duplicate-path checks, and parent-coverage summaries across delivery families.

We further validate matched-parent consistency on the released 141,811-sample test split derived from 8,407 parents. Using a transformers ASR backend (Whisper-large-v3-turbo) and a transformers speaker-embedding backend (WavLM-base-sv), we compare each delivered sample against its clean parent through transcript error-rate drift, speaker cosine similarity, duration ratio, and absolute RMS drift. Full family-wise statistics are deferred to the appendix.